\section{Detección de eventos en fútbol}

La detección de eventos en fútbol tiene como objetivo identificar cuándo ocurren acciones relevantes del partido y asignarles una interpretación semántica, como pase, tiro, falta, saque o cambio de posesión. Esta tarea es especialmente compleja porque muchas acciones no pueden reconocerse únicamente a partir de un único frame, sino que dependen de la evolución temporal de jugadores, balón y contexto de juego.

Una línea de trabajo relacionada es el \textit{action spotting}, cuyo objetivo consiste en localizar temporalmente acciones relevantes dentro de vídeos largos de partido. En fútbol, SoccerNet \citep{GiancolaSoccerNet2018} ha sido una de las principales referencias para esta tarea, al proporcionar un amplio dataset de partidos broadcast con anotaciones temporales de eventos. Posteriormente, SoccerNet-v2 \citep{DeliegeSoccerNetV22021} amplió el dataset con anotaciones más densas y nuevas tareas de comprensión de vídeo, incluyendo acciones, cambios de cámara y repeticiones.

En vídeo broadcast, muchos trabajos abordan la detección de acciones directamente a partir de la señal visual. Una línea clásica consiste en combinar extractores convolucionales con modelos temporales: primero se obtiene una representación visual de cada frame mediante una CNN y después se modela la evolución temporal de la secuencia con redes recurrentes o LSTM, como en LRCN \citep{DonahueLRCN2015}. Otra línea emplea arquitecturas espaciotemporales, como I3D \citep{CarreiraI3D2017}, capaces de capturar conjuntamente apariencia y movimiento a partir de convoluciones 3D. Estos enfoques son potentes, pero suelen requerir grandes cantidades de vídeo anotado, diversidad visual y un coste computacional considerable.

Dentro del contexto de SoccerNet, se han propuesto métodos específicos para \textit{action spotting} en fútbol broadcast. \citet{CioppaCALF2020} introducen CALF, una función de pérdida consciente del contexto temporal que permite localizar acciones alrededor de un instante anotado, evitando tratar cada frame de forma completamente aislada. Posteriormente, NetVLAD++ \citep{GiancolaNetVLAD2021} propone una agregación temporal que separa el contexto anterior y posterior a una acción, mejorando la representación de eventos en secuencias largas. Más recientemente, T-DEED \citep{XarlesTDEED2024} busca aumentar la resolución y discriminabilidad temporal de las predicciones para eventos deportivos. Este tipo de métodos resulta especialmente relevante para tareas como BAS (\textit{Ball Action Spotting}), donde se pretende detectar acciones relacionadas con el balón directamente desde el vídeo.

Frente a los enfoques basados principalmente en frames RGB, PathCRF \citep{KimPathCRF2026} plantea una alternativa basada en trayectorias estructuradas de jugadores. Como se describe con más detalle en la \hyperref[sec:pathcrf-tecnologias]{Sección~\ref*{sec:pathcrf-tecnologias}}, el modelo infiere una secuencia de estados de posesión y extrae eventos a partir de sus cambios. Su entrenamiento y evaluación se apoyan en el \textit{Sportec Open DFL Dataset}, publicado por \citet{BassekIntegratedDataset2025}, que proporciona eventos anotados y datos posicionales oficiales de partidos de la Bundesliga.

En resumen, métodos como LRCN, I3D, CALF, NetVLAD++ o T-DEED representan enfoques consolidados para detectar eventos desde vídeo broadcast. Sin embargo, PathCRF resulta especialmente interesante para pipelines que ya construyen una representación geométrica y temporal del partido, ya que permite transformar trayectorias de jugadores en eventos semánticos utilizables por etapas posteriores.